\section{Discussion}
\label{sec:discussion}

\paragraph{What the platform enables.}
The comparisons this research area needs are now configuration rather than reimplementation.
Contrasting intervention presentation forms, commit boundaries, or decision strategies under identical sensing is a matter of session templates and module selection, with counterbalancing supplied by the platform's randomised item ordering and evidence supplied by the session record.
The unbundled study that our own confound calls for, crossing the revised restore with and without the highlight cue, is exactly such a configuration.
Beyond single studies, the exported records form a corpus on which an external decision provider can be trained offline and then attached through the same seam it will be judged at, closing the loop the programme's envisaged AI-driven adaptation requires \cite{dtucompute2025rtr}.

\paragraph{Measure the quality attribute you claim.}
The over-repositioning result is the strongest argument for a design principle the platform embodies: a system should measure the property it claims to provide, so it can report its failures as readily as its successes.
Context preservation was a headline feature, yet its original geometry made small reflows worse, and nothing in ordinary operation would have surfaced that; the per-intervention residual in the session record did.
The same principle governs the module boundary, whose out-of-process cost is not asserted but recorded per proposal on a single clock.
We suggest this transfers beyond reading platforms: any interface that changes state during a continuous user task can hold the change until a natural boundary, re-anchor the user's focus across it, and record what the change cost.

\paragraph{Limitations.}
The behavioural comparison rests on four participants, drawn from the authors and acquaintances, reading one article under manual intervention; the numbers are descriptive, and the enabled condition bundles the restore with the highlight cue, so no causal attribution to the restore geometry is possible.
The revised restore is verified geometrically over the deterministic sweep but has not been re-evaluated with readers, so its behavioural benefit is unconfirmed.
The automated decision path was exercised within budget in a single advisory-mode run with a rule-based reference provider on one host, not at campaign scale or over a network, and no learned decision model has yet run against the platform.
The provider contract has one genuinely external implementation; a single outside integration is encouraging, not proof of independence.
Latency figures are loopback figures of the intended single-host deployment.
Reading material is restricted to reflowable Markdown, and our participants were general readers, not the vision-impaired or dyslexic readers who motivate the programme, so no claim here extends to those groups.

\paragraph{Future work.}
Three studies follow directly.
First, the three-condition study separating restore geometry from highlighting on the revised mechanism, which would also replace the authors-as-participants sample.
Second, a campaign-scale run of the automated decision path with a learned provider, turning the single-run latency check into a measured result.
Third, the properly powered efficacy study the platform was built to host, asking whether specific interventions improve reading rather than only measuring their immediate cost.
On the platform side, additional external providers and webcam or multimodal sources behind the existing sensing port would push the contract from author-validated toward independently validated.
